% =====================================================================
% Spec module — gemm_main_fsm (spec baseline v1.0)
% Bố cục mục theo spec_1/modules/main_fsm/latex/main.tex
% =====================================================================
\documentclass[10pt,a4paper]{article}
\newcommand{\DocHeader}{Spec module · gemm\_main\_fsm}
\newcommand{\DocVersion}{baseline v1.0 · theo đặc tả hệ thống baseline v1.0}
\newcommand{\DocFooter}{Bộ tăng tốc GEMM INT8 trên Chipyard SoC — spec baseline}
\newcommand{\DocPdfTitle}{Spec module gemm_main_fsm baseline v1.0}
% =====================================================================
% preamble.tex — gói, font, màu và định dạng chung (spec v1.0)
% Biên dịch bằng XeLaTeX. Bộ màu và cách đặt tiêu đề theo báo cáo tiến độ
% giai đoạn 1: đỏ = điều khiển/trạng thái, xanh = dữ liệu ma trận.
% =====================================================================
\usepackage[a4paper,top=2.2cm,bottom=2.2cm,left=2.1cm,right=2.1cm,headheight=15pt]{geometry}

% ---- Font: TeX Gyre (đủ dấu tiếng Việt, có sẵn trong TeX Live / Overleaf)
\usepackage{fontspec}
\setmainfont{texgyreheros}[
  Extension=.otf, UprightFont=*-regular, BoldFont=*-bold,
  ItalicFont=*-italic, BoldItalicFont=*-bolditalic]
\setsansfont{texgyreheros}[
  Extension=.otf, UprightFont=*-regular, BoldFont=*-bold,
  ItalicFont=*-italic, BoldItalicFont=*-bolditalic]
\newfontfamily\headfont{texgyreadventor}[
  Extension=.otf, UprightFont=*-regular, BoldFont=*-bold,
  ItalicFont=*-italic, BoldItalicFont=*-bolditalic]
\setmonofont{DejaVuSansMono}[
  Extension=.ttf, UprightFont=*, BoldFont=*-Bold,
  ItalicFont=*-Oblique, BoldItalicFont=*-BoldOblique, Scale=0.80]
% Khối code có chú thích tiếng Việt: DejaVu Sans Mono thiếu vài dấu, dùng TeX Gyre Cursor
\newfontfamily\vnmono{texgyrecursor}[
  NFSSFamily=vnmono, Extension=.otf, UprightFont=*-regular, BoldFont=*-bold,
  ItalicFont=*-italic, BoldItalicFont=*-bolditalic, Scale=0.9]

% ---- Ngôn ngữ
\usepackage{babel}
\babelprovide[import,main]{vietnamese}
\AtBeginDocument{%
  \renewcommand{\contentsname}{Mục lục}%
  \renewcommand{\listtablename}{Danh mục bảng}%
  \renewcommand{\listfigurename}{Danh mục hình}%
  \renewcommand{\tablename}{Bảng}%
  \renewcommand{\figurename}{Hình}}

% ---- Ký hiệu toán và Unicode
\usepackage{amsmath,amssymb}
\usepackage{newunicodechar}
\newunicodechar{→}{\ensuremath{\rightarrow}}
\newunicodechar{←}{\ensuremath{\leftarrow}}
\newunicodechar{↔}{\ensuremath{\leftrightarrow}}
\newunicodechar{⇒}{\ensuremath{\Rightarrow}}
\newunicodechar{≤}{\ensuremath{\le}}
\newunicodechar{≥}{\ensuremath{\ge}}
\newunicodechar{≠}{\ensuremath{\neq}}
\newunicodechar{≈}{\ensuremath{\approx}}
\newunicodechar{∈}{\ensuremath{\in}}
\newunicodechar{∉}{\ensuremath{\notin}}
\newunicodechar{∧}{\ensuremath{\wedge}}\newunicodechar{∨}{\ensuremath{\vee}}
\newunicodechar{⌈}{\ensuremath{\lceil}}
\newunicodechar{⌉}{\ensuremath{\rceil}}
\newunicodechar{Σ}{\ensuremath{\Sigma}}
\newunicodechar{−}{\ensuremath{-}}
\newunicodechar{×}{\ensuremath{\times}}
\newunicodechar{·}{\ensuremath{\cdot}}
\newunicodechar{²}{\ensuremath{^{2}}}
\newunicodechar{³}{\ensuremath{^{3}}}

% ---- Màu (khớp slides.html của báo cáo)
\usepackage{xcolor}
\definecolor{brand}{HTML}{BC1823}
\definecolor{brandsoft}{HTML}{F6DADC}
\definecolor{brandline}{HTML}{DF9DA2}
\definecolor{cream}{HTML}{FFF5EC}
\definecolor{card}{HTML}{FFFDF9}
\definecolor{ink}{HTML}{2B2B2B}
\definecolor{muted}{HTML}{6E5F5F}
\definecolor{data}{HTML}{1D5C7A}
\definecolor{datasoft}{HTML}{E2EDF2}
\definecolor{cpu}{HTML}{4A4040}
\definecolor{okgreen}{HTML}{2E7D4F}
\color{ink}

% ---- Tiêu đề mục: chữ đỏ, vạch đỏ bên trái như tiêu đề slide
\usepackage{titlesec}
\titleformat{\section}[block]
  {\headfont\bfseries\LARGE\color{brand}}
  {\thesection}{0.6em}{}
  [{\color{brandline}\titlerule[0.8pt]}]
\titlespacing*{\section}{0pt}{14pt}{8pt}
\titleformat{\subsection}[block]
  {\headfont\bfseries\large\color{brand}}
  {\thesubsection}{0.6em}{}
\titlespacing*{\subsection}{0pt}{12pt}{4pt}
\titleformat{\subsubsection}[block]
  {\headfont\bfseries\normalsize\color{ink}}
  {\thesubsubsection}{0.6em}{}
\titlespacing*{\subsubsection}{0pt}{8pt}{2pt}
\setcounter{secnumdepth}{3}
\setcounter{tocdepth}{2}

% ---- Bảng, danh sách, code
\usepackage{array,longtable,booktabs,multirow}
\usepackage{ragged2e}
\usepackage{enumitem}
\setlist{itemsep=1pt,topsep=3pt,leftmargin=1.4em}
\usepackage{fancyvrb}
\fvset{fontsize=\small,fontfamily=vnmono}
\setlength{\LTpre}{4pt}
\setlength{\LTpost}{6pt}
\setlength{\LTcapwidth}{\linewidth}
\renewcommand{\arraystretch}{1.18}
\newcolumntype{L}[1]{>{\RaggedRight\arraybackslash}p{\dimexpr#1\linewidth-2\tabcolsep-0.6pt\relax}}
\newcolumntype{C}[1]{>{\Centering\arraybackslash}p{\dimexpr#1\linewidth-2\tabcolsep-0.6pt\relax}}

% Bảng chuẩn: \begin{spectable}{cột}{chú thích}{nhãn}{hàng tiêu đề} ... \end{spectable}
% Hàng tiêu đề lặp lại khi bảng sang trang.
\newenvironment{spectable}[4]{%
  \begingroup\small
  \begin{longtable}{@{}#1@{}}
  \caption{#2}\label{#3}\\
  \toprule #4 \\ \midrule
  \endfirsthead
  \caption[]{#2 (tiếp theo)}\\
  \toprule #4 \\ \midrule
  \endhead
  \bottomrule
  \endlastfoot}
  {\end{longtable}\endgroup}
% Bảng yêu cầu ID | Yêu cầu
\newenvironment{reqtable}[2]{%
  \begin{spectable}{L{0.12}L{0.88}}{#1}{#2}%
  {\textbf{ID} & \textbf{Yêu cầu}}}
  {\end{spectable}}

% ---- Chú thích, đánh số theo mục
\usepackage{caption}
\captionsetup{font=small,labelfont={bf,color=brand},skip=4pt}
\counterwithin{table}{section}
\counterwithin{figure}{section}

% ---- Hình (TikZ)
\usepackage{graphicx}
\usepackage{tikz}
\usetikzlibrary{arrows.meta,positioning,fit,calc,backgrounds,shapes.geometric,decorations.pathreplacing}
\tikzset{
  >={Stealth[length=2.2mm]},
  blk/.style={draw=brand,line width=0.8pt,fill=card,rounded corners=2pt,align=center,font=\small,inner sep=4pt},
  blkd/.style={draw=data,line width=0.8pt,fill=datasoft,rounded corners=2pt,align=center,font=\small,inner sep=4pt},
  blkn/.style={draw=brand,line width=1.1pt,fill=brandsoft,rounded corners=2pt,align=center,font=\small,inner sep=4pt},
  blkc/.style={draw=cpu,line width=0.8pt,fill=white,rounded corners=2pt,align=center,font=\small,inner sep=4pt},
  ctl/.style={->,draw=brand,line width=0.8pt},
  dat/.style={->,draw=data,line width=1.3pt},
  lbl/.style={font=\scriptsize,align=center,inner sep=1pt},
}
\usepackage{float}
\usepackage{placeins}
\usepackage{needspace}
\AddToHook{cmd/section/before}{\FloatBarrier\Needspace{12\baselineskip}}
\AddToHook{cmd/subsection/before}{\Needspace{10\baselineskip}}
\AddToHook{cmd/subsection/before}{\FloatBarrier}

% ---- Hộp nội dung kiểu báo cáo
\usepackage[most]{tcolorbox}
% Hộp kết luận đầu mục (một câu trả lời trước, chi tiết sau)
\newtcolorbox{keypoint}{enhanced,breakable,colback=brandsoft!55,colframe=brand,
  boxrule=0pt,leftrule=3pt,arc=0pt,left=8pt,right=6pt,top=4pt,bottom=4pt,
  fontupper=\small}
% Hộp ghi chú
\newtcolorbox{notebox}[1][Ghi chú]{enhanced,breakable,colback=datasoft!60,colframe=data,
  boxrule=0pt,leftrule=3pt,arc=0pt,left=8pt,right=6pt,top=4pt,bottom=4pt,
  fontupper=\small,title={\headfont\bfseries\footnotesize\MakeUppercase{#1}},
  coltitle=data,colbacktitle=datasoft!60,attach title to upper={\par}}
% Bốn câu hỏi của mỗi giao tiếp (như slide "Giao tiếp 1/4")
\newcommand{\qitem}[2]{{\headfont\bfseries\footnotesize\color{brand}\MakeUppercase{#1}}\par\vspace{1pt}#2\par\vspace{5pt}}
\newtcolorbox{qcard}[1]{enhanced,breakable,colback=card,colframe=brandline,boxrule=0.6pt,
  arc=2pt,left=7pt,right=7pt,top=5pt,bottom=3pt,fontupper=\small,
  title={#1},fonttitle=\headfont\bfseries\small,coltitle=brand,colbacktitle=card,
  titlerule=0.4pt,titlerule style={brandline}}
\newcommand{\interfacecard}[5]{%
  \begin{qcard}{#1}
  \qitem{Trao đổi gì?}{#2}
  \qitem{Ai khởi tạo?}{#3}
  \qitem{Khi nào được chấp nhận?}{#4}
  \qitem{Khi nào sang bước tiếp?}{#5}
  \end{qcard}}

% Nhãn trạng thái giai đoạn
\newcommand{\stagetag}[1]{\tikz[baseline=(t.base)]\node[t,fill=brand,text=white,rounded corners=1.5pt,inner sep=1.6pt,font=\headfont\bfseries\scriptsize](t){#1};}
\tikzset{t/.style={}}

% ---- Đầu trang, chân trang
\usepackage{fancyhdr}
\pagestyle{fancy}
\fancyhf{}
\fancyhead[L]{\headfont\footnotesize\color{brand}\MakeUppercase{Spec module · \ModHeader}}
\fancyhead[R]{\headfont\footnotesize\color{muted}v1.0 · spec hệ thống v1.0, gồm mục 7}
\fancyfoot[L]{\footnotesize\color{muted}Bộ tăng tốc GEMM INT8 trên Chipyard SoC}
\fancyfoot[R]{\headfont\bfseries\footnotesize\color{brand}\thepage}
\renewcommand{\headrule}{{\color{brandline}\hrule height 0.6pt}}
\renewcommand{\footrule}{{\color{brandline}\hrule height 0.6pt}}
\renewcommand{\footrulewidth}{0.6pt}
\fancypagestyle{plain}{\fancyhf{}\renewcommand{\headrulewidth}{0pt}%
  \fancyfoot[R]{\headfont\bfseries\footnotesize\color{brand}\thepage}}

\setlength{\parindent}{0pt}
\setlength{\parskip}{4pt}
\linespread{1.08}

% ---- Lệnh riêng
\newcommand{\XN}{XN}
\newcommand{\appsection}[1]{%
  \refstepcounter{section}%
  \section*{Phụ lục \thesection: #1}%
  \addcontentsline{toc}{section}{Phụ lục \thesection: #1}}

\usepackage[hidelinks]{hyperref}
\usepackage{url}
\DeclareUrlCommand\sig{\urlstyle{tt}}
\makeatletter\g@addto@macro\UrlBreaks{\do\_\do\[}\makeatother
% Bản dùng được trong chú thích và tiêu đề (đối số di chuyển)
\DeclareRobustCommand\sigc[1]{\texttt{\detokenize{#1}}}
\hypersetup{pdftitle={Spec module gemm_systolic_array v1.0},pdfauthor={Nhóm đồ án GEMM}}

\newcommand{\eqbox}[1]{\begin{tcolorbox}[enhanced,colback=card,colframe=brandline,boxrule=0.6pt,arc=2pt,left=8pt,right=8pt,top=2pt,bottom=2pt]#1\end{tcolorbox}}

% ---- Bìa dùng chung: \makecover{tên module}{tiêu đề phụ}{dòng mô tả}{các dòng bảng}
\newcommand{\makecover}[4]{%
\thispagestyle{empty}
\begin{tikzpicture}[remember picture,overlay]
  \fill[cream] (current page.north west) rectangle (current page.south east);
\end{tikzpicture}
\vspace*{3cm}
{\headfont\bfseries\small\color{brand}\MakeUppercase{Spec module · phiên bản 1.0 · full option}\par}
\vspace{0.6cm}
\begin{tikzpicture}[baseline=(b.north)]\node[fill=brand,minimum width=0.25cm,minimum height=3.6cm,inner sep=0pt](b){};\end{tikzpicture}\hspace{0.45cm}\begin{minipage}[t]{0.9\linewidth}
{\headfont\bfseries\fontsize{28}{34}\selectfont\color{brand}#1\par}
\vspace{6pt}
{\headfont\bfseries\fontsize{18}{24}\selectfont\color{brand}#2\par}
\vspace{10pt}
{\large\color{ink}#3\par}
\vspace{4pt}
{\normalsize\color{muted}Bộ tăng tốc GEMM INT8 trên Chipyard SoC\par}
\end{minipage}
\vspace{2cm}

\begingroup\small
\begin{tabular}{@{}L{0.2}L{0.78}@{}}
\toprule
#4
\bottomrule
\end{tabular}
\endgroup
\clearpage
\tableofcontents
\clearpage}


\definecolor{brand}{HTML}{000000}
\definecolor{brandline}{HTML}{888888}
\definecolor{ink}{HTML}{000000}
\color{black}
\titleformat{\section}{\bfseries\Large}{\thesection.}{0.6em}{}
\titleformat{\subsection}{\bfseries\large}{\thesubsection.}{0.6em}{}
\titleformat{\subsubsection}{\bfseries}{\thesubsubsection.}{0.6em}{}
\counterwithout{table}{section}
\counterwithout{figure}{section}
\captionsetup{labelfont=bf,justification=raggedright,singlelinecheck=false}
\fancyhead[L]{\small Đặc tả thiết kế GEMM INT8}
\fancyhead[R]{\small\ModHeader}
\fancyfoot[L]{\small Bộ tăng tốc GEMM INT8}
\fancyfoot[R]{\small\thepage}
\hypersetup{pdftitle={Đặc tả thiết kế GEMM INT8},pdfauthor={Nhóm đồ án GEMM}}

\titleformat{\section}{\color{data}\bfseries\Large}{\thesection.}{0.6em}{}
\titleformat{\subsection}{\color{data}\large}{\thesubsection.}{0.6em}{}
\renewenvironment{spectable}[4]{%
  \begingroup\small
  \begin{longtable}{#1}
  \caption{#2}\label{#3}\\
  \hline #4 \\ \hline
  \endfirsthead
  \caption[]{#2 (tiếp theo)}\\
  \hline #4 \\ \hline
  \endhead
  \endlastfoot}
  {\end{longtable}\endgroup}
\setlength{\emergencystretch}{2em}

\begin{document}
{\Large\bfseries ĐẶC TẢ THIẾT KẾ gemm\_main\_fsm}\par
\vspace{2pt}
{\small\color{muted}Module của \sig{gemm_accel_top} · instance \sig{u_main_fsm}}\par
\tableofcontents
\clearpage

% ===================== 1 =====================
\section{Tổng quan}
\subsection{Tính năng chính}\label{sec:1}

Main FSM là module duy nhất ra lệnh cho DMA và array. Mỗi phép GEMM là một \textbf{job} lấy từ một descriptor. Ba engine độc lập đưa job qua các pha: \textbf{Đọc} (FETCH, kiểm tra, LOAD\_*), \textbf{Tính} (COMPUTE) và \textbf{Ghi} (STORE\_C). Bộ chọn bank quyết định khi nào một pha được bắt đầu, nên ba pha của ba job liên tiếp chạy chồng nhau mà không đụng cùng một bank. Module nằm trong miền \sig{clk} duy nhất (100 MHz, chu kỳ 10,0 ns).

\paragraph{Chức năng}\label{par:1-1}
\begin{itemize}
\item Mỗi descriptor là một job: FETCH → kiểm tra → LOAD\_BIAS (nếu \sig{POST_EN}) → LOAD\_A → LOAD\_B (trừ khi \sig{REUSE_B}) → COMPUTE → STORE\_C. Chuỗi dừng ở descriptor có \sig{LAST} = 1. Chuỗi descriptor là chế độ duy nhất
\item Mọi tham số của job (kích thước, địa chỉ, stride, cờ, lượng tử) lấy từ descriptor 16 word do DMA đọc về. Register File chỉ cấp địa chỉ descriptor đầu chuỗi \sig{i_cfg_desc_ptr}
\item Ba engine chạy đồng thời khi điều kiện bank và thứ tự job cho phép
\item Nội dung DRAM sau một tác vụ DONE giống hệt thực thi tuần tự từng job, với mọi thời gian trả lời của bus
\item FSM báo kết thúc tác vụ cho Register File bằng đúng một xung done hoặc một xung error; điều phối soft reset; cấp bộ đếm \sig{DESC_DONE_CNT}, thanh ghi \sig{ERR_DESC_ADDR} và tín hiệu đếm hiệu năng
\item FSM cấp cấu hình của job đang tính cho array và VPU, và bank đang tính cho scratchpad. \sig{DESC_PTR} chỉ đi tới FSM
\end{itemize}

\begin{notebox}[Quyết định baseline D-01]
Không có chế độ lập trình bằng thanh ghi: FSM không có cổng \sig{i_start_kind} và không có \sig{i_cfg_dim}, \sig{i_cfg_a_addr}, \sig{i_cfg_b_addr}, \sig{i_cfg_c_addr}, \sig{i_cfg_post_ctrl}, \sig{i_cfg_bias_addr}, \sig{i_cfg_qm0}, \sig{i_cfg_qshift}. Giao diện Register File → FSM chỉ còn \sig{i_start_pulse}, \sig{i_soft_req}, \sig{i_cfg_desc_ptr}. Mọi \sig{i_start_pulse} khởi động chuỗi: \sig{R_START} → \sig{R_FETCH} luôn luôn.
\end{notebox}

% ===================== 2 =====================
\subsection{Tham số}\label{sec:2}
\begin{spectable}{|L{0.2}|L{0.14}|L{0.66}|}{Tham số}{tab:2-1}{\textbf{Tham số} & \textbf{Mặc định} & \textbf{Ý nghĩa}}
\sig{P} & 8 & Kích thước array; giá trị hợp lệ duy nhất là 8 (D-04). Chỉ dùng cho \sig{HW_CONFIG} và assertion \\\hline

\sig{NB} & 2 & localparam cố định: hai bank mỗi buffer (D-09). Không có đường xử lý nào cho NB khác 2 \\\hline

\end{spectable}

% ===================== 3 =====================
\subsection{Tín hiệu giao tiếp}\label{sec:3}

\begin{spectable}{|L{0.26}|L{0.15}|L{0.14}|L{0.45}|}{Tín hiệu giao tiếp của \sigc{gemm_main_fsm}}{tab:3-1}{\textbf{Tên tín hiệu} & \textbf{Độ rộng} & \textbf{Hướng} & \textbf{Mô tả}}
\multicolumn{4}{@{}l}{\textit{Clock, reset, soft reset}} \\\hline

\sig{clk} & 1 & Input & Clock của khối.\newline Các thanh ghi cập nhật tại cạnh lên. \\\hline
\sig{rst_n} & 1 & Input & Reset tích cực mức thấp; được kích hoạt bất đồng bộ.\newline Cạnh nhả reset được đồng bộ với clock tại top trước khi đưa tới các module (nối vào \sig{rst_n_sync}). \\\hline

\sig{o_soft_clr} & 1 & Output & Xung xóa đồng bộ một chu kỳ gửi đến Register File, DMA, SRAM, array và VPU.\newline Chỉ phát sau khi cả hai kênh DMA đã rảnh; hủy trạng thái tác vụ hiện tại.\newline FSM không có cổng \sig{i_soft_clr}: tự xóa bằng trạng thái T\_CLR. \\\hline

\sig{o_soft_busy} & 1 & Output & Báo đang xử lý soft reset.\newline Trong khoảng này Register File báo BUSY = 1 và không nhận tác vụ mới. \\\hline

\multicolumn{4}{@{}l}{\textit{Register File → FSM}} \\\hline

\sig{i_start_pulse} & 1 & Input & Xung một chu kỳ báo Register File vừa chấp nhận một tác vụ.\newline \sig{DESC_PTR} đã chụp hợp lệ cùng chu kỳ. \\\hline

\sig{i_soft_req} & 1 & Input & Xung một chu kỳ yêu cầu Main FSM thực hiện soft reset.\newline Phát khi lệnh ghi CTRL có SOFT\_RESET = 1 và byte 0 được cho phép ghi. \\\hline

\sig{i_cfg_desc_ptr} & 32 & Input & Bản chụp đủ 32 bit của DESC\_PTR.\newline Địa chỉ descriptor đầu chuỗi.\newline Chụp khi chấp nhận tác vụ và giữ nguyên đến lần chấp nhận tiếp theo. \\\hline

\multicolumn{4}{@{}l}{\textit{FSM → Register File}} \\\hline

\sig{o_done_pulse} & 1 & Output & Xung một chu kỳ báo toàn bộ tác vụ kết thúc thành công.\newline Register File đặt DONE và xóa BUSY tại cạnh lên; không xuất hiện cùng err pulse. \\\hline

\sig{o_err_pulse} & 1 & Output & Xung một chu kỳ báo toàn bộ tác vụ kết thúc với lỗi.\newline Mã lỗi hợp lệ cùng chu kỳ; Register File đặt ERROR và xóa BUSY tại cạnh lên. \\\hline
\sig{o_err_code} & 4 & Output & Mã lỗi của tác vụ, được lấy mẫu cùng err pulse.\newline Giá trị 1: cấu hình.\newline Giá trị 2: đọc dữ liệu.\newline Giá trị 3: ghi kết quả; 4: đọc descriptor.\newline Giá trị 5 … 15 dự trữ, không bao giờ xuất hiện. \\\hline

\sig{o_perf_load} & 1 & Output & Báo pha LOAD dữ liệu, không bao gồm FETCH đang hoạt động.\newline Khi tín hiệu bằng 1 và tác vụ đang bận, bộ đếm hiệu năng tương ứng tăng tại cạnh lên. \\\hline
\sig{o_perf_comp} & 1 & Output & Báo pha COMPUTE đang hoạt động.\newline Khi tín hiệu bằng 1 và tác vụ đang bận, bộ đếm hiệu năng tương ứng tăng tại cạnh lên. \\\hline
\sig{o_perf_store} & 1 & Output & Báo pha STORE\_C đang hoạt động.\newline Khi tín hiệu bằng 1 và tác vụ đang bận, bộ đếm hiệu năng tương ứng tăng tại cạnh lên. \\\hline

\sig{o_err_desc_addr} & 32 & Output & Địa chỉ descriptor lỗi được cung cấp cho thanh ghi ERR\_DESC\_ADDR tại offset 0x44. \\\hline

\sig{o_desc_done_cnt} & 32 & Output & Số descriptor đã hoàn tất STORE\_C thành công, cung cấp cho thanh ghi DESC\_DONE\_CNT tại offset 0x48. \\\hline

\multicolumn{4}{@{}l}{\textit{Lệnh kênh đọc DMA}} \\\hline

\sig{o_rd_start} & 1 & Output & Xung một chu kỳ yêu cầu bắt đầu lệnh đọc DRAM.\newline Chỉ phát khi kênh rảnh.\newline Các tham số của lệnh được chốt tại cạnh lên của chu kỳ có o\_rd\_start = 1. \\\hline

\sig{o_rd_dst} & 2 & Output & Chọn đích của lệnh đọc.\newline Giá trị 0: buffer A.\newline Giá trị 1: buffer B.\newline Giá trị 2: buffer BIAS.\newline Giá trị 3: descriptor trả về Main FSM.\newline Giá trị được chốt cùng xung o\_rd\_start. \\\hline

\sig{o_rd_base} & 32 & Output & Địa chỉ byte của hàng đầu tiên trong DRAM cho lệnh đọc DRAM.\newline Địa chỉ phải căn 4 byte; được chốt cùng xung o\_rd\_start. \\\hline

\sig{o_rd_rows} & 6 & Output & Số hàng cần truyền trong lệnh đọc DRAM, từ 1 đến 32.\newline Giá trị được chốt cùng xung o\_rd\_start. \\\hline
\sig{o_rd_rwords} & 6 & Output & Số word 32 bit cần truyền trong mỗi hàng của lệnh đọc DRAM, từ 1 đến 32.\newline Giá trị được chốt cùng xung o\_rd\_start. \\\hline

\sig{o_rd_stride} & 16 & Output & Khoảng cách tính bằng byte giữa địa chỉ đầu hai hàng liên tiếp trong DRAM.\newline Phải là bội của 4 và không nhỏ hơn 4 lần số word mỗi hàng.\newline Được chốt cùng xung o\_rd\_start. \\\hline

\sig{o_rd_bank} & 1 & Output & Chọn bank đích của buffer A, B hoặc BIAS khi nạp dữ liệu; không dùng cho FETCH.\newline Giá trị được chốt cùng xung o\_rd\_start. \\\hline

\sig{o_rd_abort} & 1 & Output & Yêu cầu dừng lệnh đọc DRAM; đây là tín hiệu mức.\newline DMA ghi nhớ yêu cầu trong suốt lệnh, không khởi tạo burst mới và hoàn tất các giao dịch đã cam kết.\newline Abort không tự tạo lỗi. \\\hline

\sig{i_rd_busy} & 1 & Input & Báo kênh đọc DRAM đang bận.\newline Bằng 1 từ chu kỳ sau khi nhận o\_rd\_start đến hết chu kỳ báo done hoặc error; không nhận lệnh mới trong khoảng này. \\\hline
\sig{i_rd_done} & 1 & Input & Xung một chu kỳ báo lệnh đọc DRAM kết thúc mà không có lỗi AXI.\newline Không xuất hiện đồng thời với error.\newline Lệnh bị abort vẫn báo done nếu không gặp lỗi. \\\hline
\sig{i_rd_error} & 1 & Input & Xung một chu kỳ báo lệnh đọc DRAM kết thúc có lỗi.\newline Lỗi gồm phản hồi R khác OKAY hoặc RLAST sai vị trí.\newline Không xuất hiện đồng thời với done; chỉ báo sau khi giao dịch đang dở đã được xử lý. \\\hline

\sig{i_desc_valid} & 1 & Input & Báo có một word descriptor hợp lệ trong chu kỳ này.\newline Chỉ dùng với lệnh FETCH; bên nhận lấy dữ liệu khi valid = 1, không có ready. \\\hline
\sig{i_desc_idx} & 4 & Input & Chỉ số word trong descriptor, từ 0 đến 15.\newline Hợp lệ cùng chu kỳ với i\_desc\_valid. \\\hline
\sig{i_desc_data} & 32 & Input & Nội dung 32 bit của word descriptor có chỉ số i\_desc\_idx.\newline Hợp lệ cùng chu kỳ với i\_desc\_valid; dữ liệu của FETCH bị lỗi không được dùng để thực thi job. \\\hline

\multicolumn{4}{@{}l}{\textit{Lệnh kênh ghi DMA}} \\\hline

\sig{o_wr_start} & 1 & Output & Xung một chu kỳ yêu cầu bắt đầu lệnh ghi DRAM.\newline Chỉ phát khi kênh rảnh.\newline Các tham số của lệnh được chốt tại cạnh lên của chu kỳ có o\_wr\_start = 1. \\\hline

\sig{o_wr_fmt} & 1 & Output & Chọn định dạng dữ liệu trong buffer C.\newline Giá trị 0: INT32, mỗi phần tử chiếm một word.\newline Giá trị 1: INT8, bốn phần tử trong một word.\newline Giá trị được chốt cùng xung o\_wr\_start. \\\hline

\sig{o_wr_base} & 32 & Output & Địa chỉ byte của hàng đầu tiên trong DRAM cho lệnh ghi DRAM.\newline Địa chỉ phải căn 4 byte; được chốt cùng xung o\_wr\_start. \\\hline
\sig{o_wr_rows} & 6 & Output & Số hàng cần truyền trong lệnh ghi DRAM, từ 1 đến 32.\newline Giá trị được chốt cùng xung o\_wr\_start. \\\hline
\sig{o_wr_rwords} & 6 & Output & Số word 32 bit cần truyền trong mỗi hàng của lệnh ghi DRAM, từ 1 đến 32.\newline Giá trị được chốt cùng xung o\_wr\_start. \\\hline
\sig{o_wr_stride} & 16 & Output & Khoảng cách tính bằng byte giữa địa chỉ đầu hai hàng liên tiếp trong DRAM.\newline Phải là bội của 4 và không nhỏ hơn 4 lần số word mỗi hàng.\newline Được chốt cùng xung o\_wr\_start. \\\hline
\sig{o_wr_bank} & 1 & Output & Chọn bank nguồn của buffer C khi thực hiện STORE\_C.\newline Giá trị được chốt cùng xung o\_wr\_start. \\\hline

\sig{o_wr_abort} & 1 & Output & Yêu cầu dừng lệnh ghi DRAM; đây là tín hiệu mức.\newline DMA ghi nhớ yêu cầu trong suốt lệnh, không khởi tạo burst mới và hoàn tất các giao dịch đã cam kết.\newline Abort không tự tạo lỗi. \\\hline

\sig{i_wr_busy} & 1 & Input & Báo kênh ghi DRAM đang bận.\newline Bằng 1 từ chu kỳ sau khi nhận o\_wr\_start đến hết chu kỳ báo done hoặc error; không nhận lệnh mới trong khoảng này. \\\hline
\sig{i_wr_done} & 1 & Input & Xung một chu kỳ báo lệnh ghi DRAM kết thúc mà không có lỗi AXI.\newline Không xuất hiện đồng thời với error.\newline Lệnh bị abort vẫn báo done nếu không gặp lỗi. \\\hline
\sig{i_wr_error} & 1 & Input & Xung một chu kỳ báo lệnh ghi DRAM kết thúc có lỗi.\newline Lỗi xảy ra khi phản hồi B khác OKAY.\newline Không xuất hiện đồng thời với done; chỉ báo sau khi giao dịch đang dở đã được xử lý. \\\hline

\multicolumn{4}{@{}l}{\textit{Array, VPU và scratchpad phía tính}} \\\hline

\sig{o_pe_start} & 1 & Output & Xung một chu kỳ bắt đầu COMPUTE cho toàn bộ job.\newline Chỉ phát khi array rảnh; kích thước và cấu hình tính toán phải hợp lệ cùng chu kỳ. \\\hline

\sig{i_pe_done} & 1 & Input & Xung một chu kỳ báo COMPUTE đã hoàn tất, bao gồm hậu xử lý và ghi kết quả vào buffer C.\newline Không báo cho lượt tính bị soft clear hủy. \\\hline
\sig{i_pe_busy} & 1 & Input & Báo array đang tính.\newline Bằng 1 từ chu kỳ sau khi nhận pe start đến hết chu kỳ có pe done; về 0 khi soft clear hủy lượt tính. \\\hline

\sig{o_comp_dim} & 32 & Output & Kích thước job: \{8'h00, K[7:0], N[7:0], M[7:0]\}.\newline Bit [7:0]: M; [15:8]: N; [23:16]: K.\newline M từ 1 đến 32; N và K là bội của 4, từ 4 đến 32.\newline Phải hợp lệ khi bắt đầu lượt tính.\newline Bit [31:24] luôn bằng 0; giữ nguyên từ pe start đến hết pe done. \\\hline

\sig{o_comp_post_en} & 1 & Output & Cho phép hậu xử lý: 0 xuất INT32 nguyên gốc; 1 cộng bias, lượng tử hóa và xuất INT8.\newline Giữ ổn định từ pe start đến hết pe done. \\\hline
\sig{o_comp_relu_en} & 1 & Output & Cho phép chặn kết quả âm về 0 trước khi kẹp INT8.\newline Bị bỏ qua khi hậu xử lý tắt.\newline Giữ ổn định từ pe start đến hết pe done. \\\hline

\sig{o_comp_qm0} & 32 & Output & Hệ số nhân lượng tử M0, số nguyên 32 bit có dấu.\newline Giữ ổn định từ pe start đến hết pe done. \\\hline
\sig{o_comp_qshift} & 5 & Output & Lượng dịch bổ sung sh, số không dấu từ 0 đến 31.\newline Tổng lượng dịch trong công thức lượng tử là 31 + sh.\newline Giữ ổn định từ pe start đến hết pe done. \\\hline

\sig{o_comp_bank_a} & 1 & Output & Chọn bank của buffer A dành cho phía tính.\newline Giữ nguyên trong suốt COMPUTE. \\\hline
\sig{o_comp_bank_b} & 1 & Output & Chọn bank của buffer B dành cho phía tính.\newline Giữ nguyên trong suốt COMPUTE. \\\hline
\sig{o_comp_bank_bias} & 1 & Output & Chọn bank của buffer BIAS dành cho phía tính.\newline Giữ nguyên trong suốt COMPUTE. \\\hline
\sig{o_comp_bank_c} & 1 & Output & Chọn bank của buffer C dành cho phía tính.\newline Giữ nguyên trong suốt COMPUTE. \\\hline

\sig{o_comp_active} & 1 & Output & Báo COMPUTE đang hoạt động.\newline SRAM dùng tín hiệu này để kiểm tra quyền truy cập bằng assertion (AS-02 của SRAM); tín hiệu không thay thế các enable đọc/ghi. \\\hline

\end{spectable}

\paragraph{Giao tiếp}\label{par:3-2}
\begin{itemize}
\item Mọi port \sig{o_*} lấy từ thanh ghi, kể cả \sig{o_rd_abort} và \sig{o_wr_abort}; hai port này là tín hiệu mức (không phải xung), chỉ giải mã từ thanh ghi trạng thái của FSM: \sig{o_rd_abort} = T\_SOFT ∧ (engine Đọc ở R\_FWAIT ∨ R\_LWAIT), \sig{o_wr_abort} = T\_SOFT ∧ engine Ghi ở S\_RUN. Không dùng \sig{i_rd_busy}, \sig{i_wr_busy}, nên không có đường tổ hợp từ input tới hai port này
\item Lệnh DMA: FSM chỉ bật \sig{o_rd_start} khi kênh đọc rảnh (\sig{i_rd_busy} = 0 và không có lệnh nào đang chờ \sig{i_rd_done}/\sig{i_rd_error}); tương tự kênh ghi. Các port tham số chỉ cần hợp lệ ở chu kỳ \sig{o_*_start} = 1; DMA tự chốt
\item \sig{o_comp_dim}, \sig{o_comp_post_en}, \sig{o_comp_relu_en}, \sig{o_comp_qm0}, \sig{o_comp_qshift}, \sig{o_comp_bank_*} hợp lệ từ chu kỳ \sig{o_pe_start} = 1 tới hết chu kỳ \sig{i_pe_done} = 1 và không đổi trong khoảng đó
\item \sig{o_pe_start} chỉ khi \sig{i_pe_busy} = 0. Mỗi \sig{o_pe_start} có đúng một \sig{i_pe_done}
\item \sig{o_done_pulse} và \sig{o_err_pulse} không bao giờ cùng chu kỳ; mỗi \sig{i_start_pulse} có tối đa một trong hai; không có cái nào cho tác vụ bị soft reset hủy
\end{itemize}

\begin{notebox}[Quyết định baseline D-10]
\sig{o_comp_dim} rộng 32 bit như spec\_1, nội dung \{8'h00, K, N, M\}. Main FSM đã kiểm \sig{DIM[31:24]} = 0 trước khi phát nên bit [31:24] luôn bằng 0; array bỏ qua [31:24].
\end{notebox}

% ===================== 2 =====================
\section{Đặc tả thanh ghi}
Module không có thanh ghi ánh xạ MMIO. Bản đồ thanh ghi (\sig{CTRL}, \sig{STATUS}, \sig{HW_CONFIG}, \sig{PERF_*}, \sig{DESC_PTR}, \sig{ERR_DESC_ADDR}, \sig{DESC_DONE_CNT}) do \sig{gemm_register_file} sở hữu và được định nghĩa trong mục Đặc tả thanh ghi của đặc tả hệ thống. Cấu hình của từng job được nhận từ descriptor qua \sig{i_desc_*}; trạng thái nội bộ, điều kiện chốt và reset được mô tả bên dưới.

\section{Mô tả chức năng}
\subsection{Tổ chức chức năng}\label{sec:4}

\Needspace{12\baselineskip}
\begin{spectable}{|L{0.2}|L{0.17}|L{0.63}|}{Module con}{tab:4-1}{\textbf{Module} & \textbf{Instance} & \textbf{Vai trò}}
\sig{gemm_fsm_ctrl} & \sig{u_ctrl} & Trạng thái tác vụ; khởi động chuỗi; bản ghi lỗi; tín hiệu chặn; cấp quyền bắt đầu pha; done/error; soft reset; \sig{ERR_DESC_ADDR} \\\hline

\sig{gemm_job_check} & \sig{u_chk} (trong engine Đọc) & Tổ hợp: dựng job, kiểm kích thước, địa chỉ, cờ, bit dự trữ và stride theo các điều kiện trong mục Dựng và kiểm tra job; tính stride hiệu dụng. \\\hline

\sig{gemm_rd_engine} & \sig{u_rd} & Engine Đọc: FETCH, CHECK, chờ WAIT\_PREV, LOAD\_*, đẩy job vào slot \\\hline

\sig{gemm_cp_engine} & \sig{u_cp} & Engine Tính: lấy job từ slot, chạy COMPUTE, giữ job tới khi engine Ghi nhận \\\hline

\sig{gemm_wr_engine} & \sig{u_wr} & Engine Ghi: STORE\_C; \sig{DESC_DONE_CNT}; báo job cuối \\\hline

\sig{gemm_bank_sb} & \sig{u_sb} & Chọn bank theo con trỏ của từng buffer; bộ đếm chiếm dụng bank (scoreboard) \\\hline

\end{spectable}

Luồng job: engine Đọc → slot (một job) → engine Tính → (C\_HAND) → engine Ghi. Có bốn máy trạng thái, mã hóa one-hot: điều khiển tác vụ (\sig{gemm_fsm_ctrl}, 4 bit), engine Đọc (13 bit), engine Tính (3 bit), engine Ghi (2 bit). Trạng thái ban đầu sau reset là bit IDLE tương ứng.

\begin{notebox}[Quyết định baseline D-29 (phần trọng tài engine)]
Khi nhiều engine cùng muốn bắt đầu pha ở một chu kỳ, thứ tự ưu tiên cố định là Ghi > Tính > Đọc. Ba engine dùng các tài nguyên khác nhau (kênh ghi DMA, array, kênh đọc DMA) nên trọng tài chỉ quyết định thứ tự cấp \sig{go} khi \sig{want} trùng chu kỳ.
\end{notebox}

\subsubsection{Job}\label{sec:4.1}

Một job là struct \sig{gemm_pkg::job_t}. Engine Đọc giữ một job, slot giữ một job, engine Tính và engine Ghi mỗi bên giữ một job; tối đa bốn job cùng tồn tại.

\begin{spectable}{|L{0.18}|L{0.08}|L{0.74}|}{Trường của \sigc{job_t}}{tab:4-2}{\textbf{Trường} & \textbf{Rộng} & \textbf{Nội dung}}
\sig{idx} & 32 & Số thứ tự trong chuỗi, 0, 1, 2, … \\\hline

\sig{daddr} & 32 & Địa chỉ descriptor (cho \sig{ERR_DESC_ADDR}) \\\hline

\sig{m}, \sig{n}, \sig{k} & 6 × 3 & M, N, K (1 … 32) \\\hline

\sig{a}, \sig{b}, \sig{c}, \sig{bias} & 32 × 4 & Địa chỉ byte của A, B, C/Y, bias \\\hline

\sig{lda}, \sig{ldb}, \sig{ldc} & 16 × 3 & Stride hiệu dụng (byte), đã thay 0 bằng độ dài hàng \\\hline

\sig{qm0}, \sig{qsh} & 32, 5 & $M_0$, $sh$ \\\hline

\sig{post}, \sig{relu}, \sig{reuse}, \sig{waitp}, \sig{last} & 1 × 5 & Cờ \\\hline

\sig{next} & 32 & \sig{NEXT_PTR} (chỉ engine Đọc dùng) \\\hline

\sig{ba}, \sig{bb}, \sig{bbias}, \sig{bc} & 1 × 4 & Bank của A, B, BIAS, C (gán ở CHECK) \\\hline

\end{spectable}

% ===================== 5 =====================
\subsection{Dựng và kiểm tra job}\label{sec:5}

\subsubsection{Từ descriptor}

Descriptor dài 16 word = 64 byte, căn 64 byte. Engine Đọc nhận 16 word qua \sig{i_desc_*} và lưu word 0 … 10 vào bộ đệm. Word 11 (\sig{TAG}) được bỏ qua, mọi giá trị đều hợp lệ; word 12 … 15 chỉ cần OR dồn để kiểm bằng 0.

\begin{spectable}{|L{0.24}|L{0.76}|}{Job từ descriptor}{tab:5-2}{\textbf{Trường} & \textbf{Giá trị (w$n$ = word thứ $n$)}}
\sig{m}, \sig{n}, \sig{k} & w0[7:0], w0[15:8], w0[23:16]; w0[31:24] = 0 \\\hline

\sig{a}, \sig{b}, \sig{c} & w1, w2, w3 \\\hline

\sig{lda}, \sig{ldb} & w4[15:0], w4[31:16]; 0 → K, N \\\hline

\sig{ldc} & w5[15:0]; 0 → N nếu \sig{post}, ngược lại 4N; w5[31:16] dự trữ \\\hline

\sig{last}, \sig{waitp}, \sig{reuse}, \sig{post}, \sig{relu} & w6[0], w6[1], w6[2], w6[3], w6[4]; w6[31:5] dự trữ \\\hline

\sig{next}, \sig{bias}, \sig{qm0}, \sig{qsh} & w7, w8, w9, w10[4:0]; w10[31:5] dự trữ \\\hline

\sig{TAG} & w11, phần cứng không đọc \\\hline

\sig{idx}, \sig{daddr} & Bộ đếm descriptor của engine Đọc; con trỏ dùng để FETCH \\\hline

\end{spectable}

\paragraph{Kiểm tra descriptor (\sigc{gemm_job_check})}\label{par:5-3}
\begin{itemize}
\item Hợp lệ khi và chỉ khi mọi điều sau đúng: $1 ≤ M ≤ 32$; N, K ∈ \{4, 8, …, 32\}; w0[31:24] = 0; w5[31:16] = 0; w6[31:5] = 0; w10[31:5] = 0; w12 … w15 = 0
\item \sig{a}[1:0] = 0 và \sig{c}[1:0] = 0; \sig{b}[1:0] = 0 nếu \sig{reuse} = 0; \sig{bias}[1:0] = 0 nếu \sig{post} = 1
\item Mỗi stride thô (w4[15:0], w4[31:16], w5[15:0]) bằng 0, hoặc là bội của 4 và không nhỏ hơn độ dài hàng: K cho A; N cho B; 4N cho C khi \sig{post} = 0, N khi \sig{post} = 1
\item \sig{relu} = 1 kéo theo \sig{post} = 1
\item \sig{reuse} = 1 chỉ khi \sig{idx} > 0 và (K, N) bằng (K, N) của job trước trong chuỗi (thanh ghi \sig{prev_kn} của engine Đọc)
\item \sig{last} = 0 kéo theo \sig{next}[5:0] = 0. Vi phạm được tính là lỗi của chính descriptor này
\item \sig{DESC_PTR}[5:0] ≠ 0 được kiểm ở \sig{gemm_fsm_ctrl} khi bắt đầu chuỗi: lỗi mã 1 với f = 0, \sig{ERR_DESC_ADDR} = \sig{DESC_PTR}, không truy cập bus và không truy cập scratchpad
\end{itemize}

\begin{notebox}[Quyết định baseline D-24 và D-32]
Phần cứng chỉ kiểm các điều kiện trên (D-32, gộp cả hai nguồn). Phần cứng không kiểm vùng nhớ có nằm trong DRAM hay không và không kiểm chồng lấn giữa vùng ghi C/Y với vùng A, B, bias; vi phạm cho kết quả không xác định, lỗi bus báo bằng mã 2, 3 hoặc 4. N, K bội của 4 do phần mềm đệm 0 (D-26); M không chia hết cho 8 do phần cứng xử lý trong tile biên.
\end{notebox}

\paragraph{Mã lỗi}\label{par:5-5}
\begin{spectable}{|L{0.1}|L{0.28}|L{0.62}|}{Mã lỗi \sigc{ERR_CODE}}{tab:5-4}{\textbf{Mã} & \textbf{Nguồn} & \textbf{Điều kiện}}
0 & Không lỗi & Không có lỗi \\\hline
1 & CHECK, \sig{gemm_fsm_ctrl} & Descriptor không hợp lệ, hoặc \sig{DESC_PTR}[5:0] ≠ 0 \\\hline
2 & LOAD\_A / LOAD\_B / LOAD\_BIAS & \sig{i_rd_error} khi nạp dữ liệu \\\hline
3 & STORE\_C & \sig{i_wr_error} \\\hline
4 & FETCH & \sig{i_rd_error} khi đọc descriptor \\\hline
5 … 15 & Dự trữ & Không bao giờ xuất hiện \\\hline
\end{spectable}

\begin{notebox}[Quyết định baseline D-02]
Mã lỗi 0 … 4. Mã 5 … 15 dự trữ. Spec\_2 định nghĩa thêm mã 5 (“tính năng không có”); baseline bỏ vì chỉ có một chế độ (D-01).
\end{notebox}

% ===================== 6 =====================
\subsection{Bank và cấp quyền bắt đầu pha}\label{sec:6}

\subsubsection{Chọn bank}\label{sec:6.1}

\paragraph{\sigc{gemm_bank_sb}: chọn bank}\label{par:6-1}
\begin{itemize}
\item Thanh ghi con trỏ \sig{nxt_a}, \sig{nxt_c}, \sig{nxt_bias}, \sig{nxt_b}, \sig{prev_b}: về 0 khi chấp nhận tác vụ, khi \sig{rst_n_sync} và khi \sig{soft_clr}
\item Ở chu kỳ CHECK của job hợp lệ: \sig{ba = nxt_a}, \sig{bc = nxt_c}; \sig{bbias = nxt_bias} nếu \sig{post} (ngược lại 0); \sig{bb = prev_b} nếu \sig{reuse}, ngược lại \sig{bb = nxt_b} và \sig{prev_b <= nxt_b}
\item Cùng cạnh clock đó đảo \sig{nxt_a}, \sig{nxt_c}; đảo \sig{nxt_bias} nếu \sig{post}; đảo \sig{nxt_b} nếu \sig{reuse} = 0
\end{itemize}

\begin{notebox}[Quyết định baseline D-29 (phần chọn bank)]
A và C đảo bank sau mỗi descriptor; B đảo sau mỗi lần có LOAD\_B (\sig{REUSE_B} = 1 dùng lại bank của descriptor trước); BIAS đảo sau mỗi descriptor có \sig{POST_EN} = 1. Mỗi buffer có đúng hai bank (D-09).
\end{notebox}

\subsubsection{Bộ đếm chiếm dụng}\label{sec:6.2}
\begin{spectable}{|L{0.17}|L{0.12}|L{0.36}|L{0.35}|}{Bộ đếm của \sigc{gemm_bank_sb}}{tab:6-2}{\textbf{Bộ đếm} & \textbf{Rộng} & \textbf{Tăng khi} & \textbf{Giảm khi}}
\sig{rcnt[A][b]} & 2 & Engine Đọc bắt đầu LOAD\_A vào bank b & COMPUTE của job có \sig{ba} = b xong (\sig{i_pe_done}) \\\hline

\sig{rcnt[B][b]} & 2 & LOAD\_B vào bank b, hoặc job \sig{reuse} = 1 có \sig{bb} = b đi qua bước B & COMPUTE của job có \sig{bb} = b xong \\\hline

\sig{rcnt[BIAS][b]} & 2 & LOAD\_BIAS vào bank b & COMPUTE của job \sig{post} = 1 có \sig{bbias} = b xong \\\hline

\sig{wcnt[C][b]} & 1 & COMPUTE của job có \sig{bc} = b bắt đầu & STORE\_C của job đó xong OKAY (\sig{i_wr_done}) \\\hline

\end{spectable}

\paragraph{Điều kiện bank}\label{par:6-3}
\begin{itemize}
\item LOAD\_X vào bank b chỉ được bắt đầu khi \sig{rcnt[X][b]} = 0: mọi job đọc bank đó đã tính xong. Job \sig{reuse} = 1 không chờ (dữ liệu B của bank đó vẫn còn, vì chưa có LOAD\_B nào vào bank ấy kể từ job trước)
\item COMPUTE vào bank C b chỉ được bắt đầu khi \sig{wcnt[C][b]} = 0: STORE\_C cuối cùng đọc bank đó đã xong
\item Tăng và giảm cùng chu kỳ trên cùng một bộ đếm được cộng ròng: +1 − 1 = giữ nguyên. Trường hợp duy nhất có thể xảy ra là bước R\_RB tăng \sig{rcnt[B][bb]} đúng chu kỳ \sig{i_pe_done} giảm cùng bộ đếm đó (job đang tính dùng cùng bank B). LOAD\_X chỉ tăng khi \sig{rcnt[X][b]} = 0 và COMPUTE chỉ tăng \sig{wcnt[C][b]} khi bộ đếm = 0, nên hai trường hợp này không trùng với lần giảm trên cùng bộ đếm
\item Bộ đếm không bao giờ tràn hay âm (assertion). \sig{rcnt} tối đa 3 vì tối đa ba job đã gán bank mà chưa tính xong (engine Đọc, slot, engine Tính)
\item Khi tác vụ kết thúc do lỗi, bộ đếm có thể khác 0; chúng được xóa khi chấp nhận tác vụ kế tiếp hoặc khi \sig{soft_clr}
\item Các điều kiện trên là điều kiện cần để bắt đầu pha; phần cứng được phép bắt đầu muộn hơn, không bao giờ sớm hơn
\end{itemize}

\subsubsection{Cấp quyền bắt đầu pha}\label{sec:6.3}
Mỗi engine tính tín hiệu \sig{want} từ trạng thái đầu chu kỳ: \sig{want_s} (Ghi), \sig{want_c} (Tính), \sig{want_r} (Đọc: FETCH hoặc một LOAD). \sig{gemm_fsm_ctrl} trả \sig{go}.

\[
go\_x = want\_x \wedge \neg inh
\]
\sig{inh} = (ctrl không ở T\_RUN) ∨ \sig{i_soft_req}. Thành phần \sig{i_soft_req} là bắt buộc: ở chu kỳ \sig{i_soft_req} đến, ctrl vẫn còn ở T\_RUN; nếu chỉ xét trạng thái thì \sig{go} vẫn được cấp và \sig{o_rd_start}, \sig{o_wr_start} hoặc \sig{o_pe_start} sẽ lên ở chu kỳ kế tiếp, khi ctrl đã ở T\_SOFT (vi phạm AS-12). Ba engine được cấp quyền độc lập; thứ tự job do bộ đếm bank giữ. Điều kiện bắt đầu pha của job d:
\begin{itemize}
\item LOAD vào một bank bắt đầu khi COMPUTE cuối cùng đọc bank đó đã xong
\item COMPUTE(d) bắt đầu khi mọi LOAD của d xong, COMPUTE(d−1) xong, và STORE\_C cuối cùng đọc bank C của d xong
\item STORE\_C(d) bắt đầu khi COMPUTE(d) và STORE\_C(d−1) xong
\item \sig{WAIT_PREV} = 1 chặn các pha LOAD của descriptor (chờ \sig{desc_done_cnt} bằng \sig{idx}); không chặn FETCH và kiểm tra
\end{itemize}

\paragraph{Chặn sau lỗi}\label{par:6-4}
\begin{itemize}
\item \sig{gemm_fsm_ctrl} cấp \sig{blk_v}, \sig{blk_idx} tổ hợp: bản ghi lỗi hiện có, gộp với lỗi đến trong chính chu kỳ này (\sig{i_rd_error}, \sig{i_wr_error}, kết quả CHECK sai), lấy chỉ số nhỏ nhất
\item Mọi \sig{want} của job có \sig{idx} ≥ \sig{blk_idx} bằng 0 khi \sig{blk_v} = 1. Nhờ gộp lỗi cùng chu kỳ, không pha nào của descriptor ≥ f bắt đầu, kể cả ở đúng chu kỳ lỗi được báo
\item Pha đã bắt đầu thì chạy tới hết (không abort khi lỗi). Kết quả của chúng không được ghi ra DRAM vì STORE\_C của job ≥ f không bao giờ bắt đầu
\end{itemize}

% ===================== 7 =====================
\subsection{Engine}\label{sec:7}

\subsubsection{Lệnh DMA}\label{sec:7.1}
\begin{spectable}{|L{0.13}|L{0.12}|L{0.13}|L{0.1}|L{0.16}|L{0.14}|L{0.12}|}{Tham số lệnh DMA của mỗi pha}{tab:7-1}{\textbf{Pha} & \textbf{dst / fmt} & \textbf{base} & \textbf{rows} & \textbf{rwords} & \textbf{stride} & \textbf{bank}}
FETCH & 3 (desc) & con trỏ & 1 & 16 & 64 & 0 \\\hline

LOAD\_BIAS & 2 (BIAS) & \sig{bias} & 1 & N & 4N & \sig{bbias} \\\hline

LOAD\_A & 0 (A) & \sig{a} & M & K/4 & \sig{lda} & \sig{ba} \\\hline

LOAD\_B & 1 (B) & \sig{b} & K & N/4 & \sig{ldb} & \sig{bb} \\\hline

STORE\_C & \sig{post} & \sig{c} & M & \sig{post} ? N/4 : N & \sig{ldc} & \sig{bc} \\\hline

\end{spectable}

FETCH là một burst INCR 16 beat; descriptor không đi vào scratchpad mà DMA trả về qua \sig{i_desc_valid}/\sig{i_desc_idx}/\sig{i_desc_data}.

\subsubsection{Engine Đọc}\label{sec:7.2}
Mã hóa one-hot 13 bit, thứ tự bit: R\_IDLE, R\_START, R\_FETCH, R\_FWAIT, R\_CHECK, R\_WAITP, R\_LBIAS, R\_LA, R\_LB, R\_RB, R\_LWAIT, R\_PUSH, R\_STOP.

\begin{spectable}{|L{0.11}|L{0.3}|L{0.3}|L{0.29}|}{Trạng thái của \sigc{gemm_rd_engine}}{tab:7-2}{\textbf{Trạng thái} & \textbf{Làm gì} & \textbf{Ra khỏi trạng thái khi} & \textbf{Sang}}
R\_IDLE & Chờ & ctrl khởi động tác vụ (\sig{i_start_pulse} ở T\_IDLE) & R\_START; R\_STOP nếu \sig{DESC_PTR}[5:0] ≠ 0 \\\hline

R\_START & Con trỏ ← \sig{DESC_PTR}, idx ← 0 & Sau đúng 1 chu kỳ & R\_FETCH (luôn luôn) \\\hline

R\_FETCH & \sig{want_r} = 1; khi \sig{go_r}: phát FETCH & \sig{go_r} = 1 & R\_FWAIT \\\hline

R\_FWAIT & Nhận word descriptor; chờ \sig{i_rd_done} / \sig{i_rd_error} & \sig{i_rd_done} = 1, hoặc \sig{i_rd_error} = 1 & R\_CHECK; R\_STOP (lỗi mã 4) \\\hline

R\_CHECK & 1 chu kỳ: \sig{gemm_job_check}; nếu hợp lệ, gán bank & Sau đúng 1 chu kỳ & R\_WAITP; R\_STOP (lỗi mã 1) \\\hline

R\_WAITP & Nếu \sig{waitp} = 1: chờ \sig{desc_done_cnt == idx} & \sig{waitp} = 0, hoặc \sig{desc_done_cnt} = \sig{idx} & R\_LBIAS nếu \sig{post}, ngược lại R\_LA \\\hline

R\_LBIAS & \sig{want_r} = (\sig{rcnt[BIAS][bbias]} = 0). Khi \sig{go_r}: \sig{rcnt[BIAS][bbias]++}, phát LOAD\_BIAS & \sig{go_r} = 1 & R\_LWAIT \\\hline

R\_LA & \sig{want_r} = (\sig{rcnt[A][ba]} = 0). Khi \sig{go_r}: \sig{rcnt[A][ba]++}, phát LOAD\_A & \sig{go_r} = 1 & R\_LWAIT \\\hline

R\_LB & \sig{want_r} = (\sig{rcnt[B][bb]} = 0). Khi \sig{go_r}: \sig{rcnt[B][bb]++}, phát LOAD\_B & \sig{go_r} = 1 & R\_LWAIT \\\hline

R\_RB & Job \sig{reuse}: \sig{rcnt[B][bb]++}, không phát lệnh & Sau đúng 1 chu kỳ & R\_PUSH \\\hline

R\_LWAIT & Chờ \sig{i_rd_done} / \sig{i_rd_error} & \sig{i_rd_done} = 1, hoặc \sig{i_rd_error} = 1 & Thứ tự BIAS → A → (B hoặc RB) → PUSH; R\_STOP (lỗi mã 2) \\\hline

R\_PUSH & Khi slot trống: slot ← job. Nếu \sig{last} = 0: \sig{prev_kn} ← (K, N), con trỏ ← \sig{next}, idx++ & Slot trống & R\_FETCH nếu \sig{last} = 0; R\_STOP nếu \sig{last} = 1 \\\hline

R\_STOP & Không làm gì tới hết tác vụ & Reset, T\_CLR, hoặc ctrl chuyển T\_RUN → T\_IDLE (DONE hoặc ERROR, D-34) & R\_IDLE \\\hline

\end{spectable}

\paragraph{Engine Đọc}\label{par:7-3}
\begin{itemize}
\item Ở mọi trạng thái chờ (R\_FETCH, R\_WAITP, R\_LBIAS, R\_LA, R\_LB, R\_PUSH): nếu job bị chặn thì sang R\_STOP, không đẩy job vào slot
\item Engine Đọc làm xong mọi LOAD của job d rồi mới FETCH job d+1. \sig{WAIT_PREV} không chặn FETCH và CHECK
\item Mỗi descriptor được FETCH đúng một lần, theo đúng thứ tự chuỗi (D-23)
\item Bộ đệm descriptor ghi word \sig{i_desc_idx} khi \sig{i_desc_valid}. Dữ liệu của FETCH bị lỗi không được dùng
\item \sig{o_perf_load} = 1 khi ở R\_LWAIT (không tính FETCH)
\item Trình tự job thường: R\_FETCH → R\_FWAIT → R\_CHECK → R\_WAITP → (R\_LBIAS) → R\_LA → R\_LB (hoặc R\_RB) → R\_PUSH → R\_FETCH; mỗi LOAD đi qua R\_LWAIT
\end{itemize}

\subsubsection{Engine Tính và engine Ghi}\label{sec:7.3}
Engine Tính mã hóa one-hot 3 bit (C\_IDLE, C\_RUN, C\_HAND); engine Ghi one-hot 2 bit (S\_IDLE, S\_RUN).

\begin{spectable}{|L{0.11}|L{0.31}|L{0.29}|L{0.29}|}{Trạng thái của \sigc{gemm_cp_engine} và \sigc{gemm_wr_engine}}{tab:7-4}{\textbf{Trạng thái} & \textbf{Làm gì} & \textbf{Ra khỏi trạng thái khi} & \textbf{Sang}}
C\_IDLE & \sig{want_c} = slot có job ∧ \sig{wcnt[C][slot.bc]} = 0. Khi \sig{go_c}: job ← slot, slot trống, \sig{wcnt++}, \sig{o_pe_start} = 1 ở chu kỳ kế tiếp & \sig{go_c} = 1 & C\_RUN \\\hline

C\_RUN & \sig{o_comp_active} = \sig{o_perf_comp} = 1. Khi \sig{i_pe_done}: giảm \sig{rcnt[A][ba]}, \sig{rcnt[B][bb]}, và \sig{rcnt[BIAS][bbias]} nếu \sig{post} & \sig{i_pe_done} = 1 & C\_HAND \\\hline

C\_HAND & Giữ job (và bank C) cho engine Ghi & Engine Ghi nhận job (\sig{go_s} = 1) & C\_IDLE \\\hline

S\_IDLE & \sig{want_s} = engine Tính ở C\_HAND. Khi \sig{go_s}: job ← job của engine Tính, phát STORE\_C ở chu kỳ kế tiếp & \sig{go_s} = 1 & S\_RUN \\\hline

S\_RUN & \sig{o_perf_store} = 1. \sig{i_wr_done}: \sig{wcnt[C][bc]--}; \sig{desc_done_cnt++}; nếu \sig{last} = 1, báo job cuối xong. \sig{i_wr_error}: lỗi mã 3 & \sig{i_wr_done} = 1, hoặc \sig{i_wr_error} = 1 & S\_IDLE \\\hline

\end{spectable}

\begin{notebox}[Vì sao engine Tính giữ job ở C\_HAND không làm chậm]
COMPUTE(d+1) dùng bank C của job d−1, nên theo điều kiện bắt đầu COMPUTE (mọi LOAD của d+1 xong, COMPUTE(d) xong, STORE\_C cuối cùng đọc bank C của d+1 xong) chỉ bắt đầu được khi STORE\_C(d−1) xong. Đúng lúc đó engine Ghi rảnh và nhận job d, nên engine Tính về C\_IDLE và có thể bắt đầu d+1 ngay chu kỳ sau.
\end{notebox}

% ===================== 8 =====================
\subsection{Điều khiển tác vụ}\label{sec:8}

Mã hóa one-hot 4 bit, thứ tự bit: T\_IDLE, T\_RUN, T\_SOFT, T\_CLR; trạng thái reset là T\_IDLE.

\begin{spectable}{|L{0.11}|L{0.33}|L{0.28}|L{0.28}|}{Trạng thái của \sigc{gemm_fsm_ctrl}}{tab:8-1}{\textbf{Trạng thái} & \textbf{Làm gì} & \textbf{Ra khỏi trạng thái khi} & \textbf{Sang}}
T\_IDLE & Chờ. \sig{i_start_pulse}: xóa bản ghi lỗi, \sig{desc_done_cnt}, \sig{err_desc_addr}, con trỏ bank, bộ đếm; khởi động engine Đọc & \sig{i_start_pulse} = 1; hoặc \sig{i_soft_req} = 1 & T\_RUN; T\_CLR \\\hline

T\_RUN & Engine chạy. Kiểm điều kiện kết thúc. Ở cạnh clock chuyển sang T\_IDLE: dọn engine (D-34): engine Đọc → R\_IDLE, engine Tính → C\_IDLE, engine Ghi → S\_IDLE, \sig{slot_v} ← 0 & \sig{i_soft_req} = 1 (ưu tiên); hoặc điều kiện DONE hoặc ERROR đúng & T\_SOFT; T\_IDLE \\\hline

T\_SOFT & \sig{inh} = 1; \sig{o_rd_abort} = 1 khi engine Đọc ở R\_FWAIT hoặc R\_LWAIT, \sig{o_wr_abort} = 1 khi engine Ghi ở S\_RUN (tức kênh đó đang có lệnh); không báo done/err & Không engine nào chờ lệnh DMA & T\_CLR \\\hline

T\_CLR & \sig{o_soft_clr} = 1 trong 1 chu kỳ; mọi thanh ghi của FSM về giá trị reset & Sau đúng 1 chu kỳ & T\_IDLE \\\hline

\end{spectable}

\paragraph{Kết thúc tác vụ}\label{par:8-3}
\begin{itemize}
\item Bản ghi lỗi (\sig{err_v}, \sig{err_idx}, \sig{err_code}, \sig{err_addr}) chỉ được thay khi lỗi mới có idx nhỏ hơn. \sig{err_addr} = \sig{daddr} của job lỗi; = con trỏ đang FETCH với lỗi mã 4; = \sig{DESC_PTR} khi \sig{DESC_PTR}[5:0] ≠ 0 (lỗi mã 1 với f = 0)
\item DONE: ở T\_RUN, engine Ghi báo job có \sig{last} = 1 xong OKAY và \sig{err_v} = 0 → \sig{o_done_pulse} ở chu kỳ kế tiếp. Chuỗi kết thúc ở descriptor đầu tiên có \sig{LAST} = 1
\item ERROR: ở T\_RUN, \sig{err_v} = 1 và hệ yên: engine Đọc ở R\_STOP và không chờ DMA; engine Tính không ở C\_RUN; slot trống hoặc có idx ≥ f; engine Tính không giữ ở C\_HAND một job idx < f; engine Ghi ở S\_IDLE → \sig{o_err_pulse} với \sig{o_err_code} = \sig{err_code} ở chu kỳ kế tiếp; \sig{o_err_desc_addr} ← \sig{err_addr}
\item Khi đó \sig{o_desc_done_cnt} = f tự nhiên, vì mọi job < f đã ghi xong và không job ≥ f nào được ghi
\item \sig{DESC_DONE_CNT} tăng 1 cho mỗi STORE\_C nhận phản hồi B cuối là OKAY
\item Dọn engine khi kết thúc (D-34): ở cạnh clock ctrl rời T\_RUN về T\_IDLE (cùng cạnh phát \sig{o_done_pulse} hoặc \sig{o_err_pulse}), engine Đọc về R\_IDLE, engine Tính về C\_IDLE, engine Ghi về S\_IDLE và \sig{slot_v} ← 0. Job idx ≥ f còn trong slot hoặc ở C\_HAND bị bỏ, không bao giờ được COMPUTE hay STORE\_C. Việc dọn an toàn vì điều kiện DONE/ERROR đã bảo đảm không kênh DMA nào bận và engine Tính không ở C\_RUN. Không có bước này, engine Đọc kẹt ở R\_STOP và tác vụ kế tiếp không bao giờ FETCH (BUSY = 1 mãi), còn job cũ sót lại sẽ bị tính và ghi vào \sig{C_ADDR} của tác vụ trước
\end{itemize}

\begin{notebox}[Quyết định baseline D-23 và D-11]
\sig{o_done_pulse} và \sig{o_err_pulse} không bao giờ cùng một chu kỳ. Khi \sig{DESC_PTR}[5:0] ≠ 0, ctrl ghi bản ghi lỗi (mã 1, f = 0) ngay chu kỳ \sig{i_start_pulse}, engine Đọc đi thẳng R\_IDLE → R\_STOP, không có lệnh DMA nào được phát, và \sig{o_err_pulse} xuất hiện ở chu kỳ t + 2 tính từ \sig{i_start_pulse} ở chu kỳ t.
\end{notebox}

\paragraph{Soft reset}\label{par:8-4}
\begin{itemize}
\item \sig{o_soft_busy} = 1 ở T\_SOFT và T\_CLR; trong cả chuỗi Register File đọc STATUS ra 0x0000\_0002, START và SOFT\_RESET ghi lặp bị bỏ qua
\item Ở T\_SOFT: không pha mới nào bắt đầu, không lệnh DMA mới; lệnh DMA đang chạy kết thúc bằng done hoặc error như thường (burst đã cam kết chạy hết); engine vẫn nhận các xung đó để rời trạng thái chờ, nhưng không cập nhật bản ghi lỗi hay báo done
\item COMPUTE đang chạy không được chờ: \sig{o_soft_clr} hủy array và VPU
\item \sig{i_soft_req} ở T\_IDLE: T\_CLR ở chu kỳ kế tiếp. \sig{i_soft_req} ở T\_SOFT hoặc T\_CLR: bỏ qua
\item \sig{i_soft_req} ở T\_RUN thắng điều kiện DONE/ERROR đúng trong cùng chu kỳ: ctrl sang T\_SOFT, không phát \sig{o_done_pulse} hay \sig{o_err_pulse} (tác vụ kết thúc bằng soft reset, không DONE, không ERROR)
\item Nếu \sig{i_soft_req} đến cùng chu kỳ với \sig{i_rd_done} hoặc \sig{i_wr_done}, kênh đó coi như đã rảnh
\item Sau \sig{o_soft_clr}: STATUS = 0, không báo DONE, không báo ERROR; \sig{DESC_PTR} và các thanh ghi RW giữ nguyên; nội dung scratchpad không xác định. Abort không phải lỗi
\end{itemize}

\begin{notebox}[Quyết định baseline D-33]
Trình tự: \sig{CTRL}[1] = 1 → \sig{o_soft_req} (thắng START trong cùng lệnh ghi) → T\_SOFT (T\_RUN) hoặc T\_CLR (T\_IDLE) → \sig{o_soft_clr} đúng một chu kỳ tới Register File, DMA, SRAM, array, VPU → T\_IDLE. FSM tự reset, không có \sig{i_soft_clr}.
\end{notebox}

% ===================== 9 =====================
\subsection{Timing}\label{sec:9}

\begin{spectable}{|L{0.58}|L{0.42}|}{Độ trễ cố định của FSM}{tab:9-1}{\textbf{Sự kiện} & \textbf{Chu kỳ}}
\sig{i_start_pulse} → engine Đọc ở R\_START & +1 (10 ns) \\\hline

R\_CHECK & 1 chu kỳ (10 ns) \\\hline

\sig{go_*} = 1 → \sig{o_rd_start} / \sig{o_wr_start} / \sig{o_pe_start} = 1 & +1 (10 ns) \\\hline

\sig{i_*_done} → trạng thái kế tiếp của engine & +1 (10 ns) \\\hline

Điều kiện DONE / ERROR → \sig{o_done_pulse} / \sig{o_err_pulse} & +1 (10 ns) \\\hline

\sig{i_soft_req} ở T\_IDLE → \sig{o_soft_clr} & +1 (10 ns) \\\hline

\sig{i_start_pulse} → \sig{o_err_pulse} khi \sig{DESC_PTR}[5:0] ≠ 0 & +2 (20 ns) \\\hline

\end{spectable}

Bảng dưới là số liệu của các module khác mà FSM phải tính đến khi lập lịch. Một chu kỳ ở 100 MHz là 10 ns.

\begin{spectable}{|L{0.6}|L{0.4}|}{Thời gian các pha cho M = N = K = 32, \sigc{POST_EN} = 0, stride = 0}{tab:9-2}{\textbf{Pha} & \textbf{Thời gian}}
FETCH & 16 beat \\\hline
LOAD\_A & 32 × 8 = 256 beat \\\hline
LOAD\_B & 32 × 8 = 256 beat \\\hline
COMPUTE & 534 chu kỳ = 5,34 µs \\\hline
STORE\_C & 32 × 32 = 1024 beat \\\hline
Đỉnh array (tham khảo) & 64 MAC/chu kỳ × 100 MHz = 6,4 GMAC/s \\\hline
\end{spectable}

\begin{notebox}[Quyết định baseline D-06 và D-25]
COMPUTE\_CYCLES = (NT − 1)·max(K, P) + K + P + r\_last + LAT + 3, với r\_last = (số hàng hợp lệ của tile cuối) − 1 và LAT = 4. Với M = N = K = 32: 15·32 + 32 + 8 + 7 + 4 + 3 = 534 chu kỳ. Số liệu thời gian nạp và ghi do người soạn baseline tính từ hai nguồn, ở mức 1 beat/chu kỳ và bỏ qua overhead burst. FETCH + nạp + ghi = 1552 beat ≈ 15,52 µs cho một descriptor chạy đơn lẻ, lớn hơn COMPUTE. Kênh đọc và kênh ghi AXI độc lập, nên khi chạy chồng với ping-pong thời gian mỗi descriptor tiến tới max(FETCH + nạp, tính, ghi) = max(528, 534, 1024) chu kỳ ≈ thời gian ghi: tác vụ 32³ bị giới hạn bởi kênh ghi AXI 32 bit, không bởi array.
\end{notebox}

\begin{figure}[H]
\centering
\begin{wave}{9}{-10.4}
\sname{0}{Read engine} \bus{0}{2}{0}{LOAD(d)} \bus{2}{3}{0}{FETCH(d+1)} \bus{3}{5}{0}{LOAD(d+1)} \bus{5}{6}{0}{FETCH(d+2)} \bus{6}{9}{0}{LOAD(d+2)}
\sname{-1.3}{Compute engine} \bus{0}{2}{-1.3}{COMPUTE(d-1)} \bus{2}{5}{-1.3}{COMPUTE(d)} \bus{5}{9}{-1.3}{COMPUTE(d+1)}
\sname{-2.6}{Write engine} \bus{0}{2}{-2.6}{STORE\_C(d-2)} \bus{2}{5}{-2.6}{STORE\_C(d-1)} \bus{5}{9}{-2.6}{STORE\_C(d)}
\sname{-3.9}{o\_pe\_start} \pulse{0}{-4.1} \lo{1}{2}{-4.1} \pulse{2}{-4.1} \lo{3}{5}{-4.1} \pulse{5}{-4.1} \lo{6}{9}{-4.1}
\sname{-5.2}{o\_wr\_start} \pulse{0}{-5.4} \lo{1}{2}{-5.4} \pulse{2}{-5.4} \lo{3}{5}{-5.4} \pulse{5}{-5.4} \lo{6}{9}{-5.4}
\sname{-6.5}{o\_comp\_bank\_a / c} \bus{0}{2}{-6.5}{0} \bus{2}{5}{-6.5}{1} \bus{5}{9}{-6.5}{0}
\sname{-7.8}{o\_wr\_bank} \bus{0}{2}{-7.8}{1} \bus{2}{5}{-7.8}{0} \bus{5}{9}{-7.8}{1}
\sname{-9.1}{o\_comp\_active} \hi{0}{9}{-9.3}
\end{wave}
\caption{Chồng pha của chuỗi descriptor (trục thời gian không theo tỷ lệ; mỗi cột là một mốc, các pha thực có độ dài khác nhau). Ở mốc 2 … 5, FETCH(d+1) và LOAD(d+1) chạy cùng COMPUTE(d) trên bank \sigc{o_comp_bank_a} = \sigc{o_comp_bank_c} = 1 và STORE\_C(d−1) đọc bank C = 0. LOAD(d+1) ghi vào bank A của d−1 (đã tính xong ở mốc 2). COMPUTE(d+1) bắt đầu ở mốc 5, đúng khi STORE\_C(d−1) xong và bank C = 0 được giải phóng; STORE\_C(d) bắt đầu cùng lúc khi COMPUTE(d) xong. Vì trục là mốc, các xung vẽ cùng cột không có nghĩa là cùng chu kỳ: ở mức chu kỳ, \sigc{o_wr_start} và \sigc{o_pe_start} của cùng một mốc lệch nhau ít nhất 1 chu kỳ (engine Tính chỉ rời C\_HAND khi engine Ghi nhận job bằng \sigc{go_s}, rồi mới có \sigc{go_c} cho job kế tiếp), và \sigc{o_comp_active} có khoảng 0 ngắn giữa hai COMPUTE (C\_HAND, C\_IDLE); xem bảng trạng thái của engine Tính và engine Ghi}\label{fig:9-1}
\end{figure}

% ===================== 10 =====================
\subsection{Reset, hiện thực, kiểm chứng}\label{sec:10}

\paragraph{Reset}\label{par:10-1}
\begin{itemize}
\item \sig{rst_n_sync} = 0 hoặc T\_CLR: mọi engine về trạng thái IDLE, slot trống, bản ghi lỗi, con trỏ bank, bộ đếm, \sig{desc_done_cnt}, \sig{err_desc_addr}, mọi xung đầu ra về 0. Trường dữ liệu của job (\sig{job_t}, bộ đệm descriptor) không reset và dùng \sig{always_ff @(posedge clk)}
\item Flip-flop trạng thái dùng \sig{always_ff @(posedge clk or negedge rst_n_sync)}; T\_CLR xóa đồng bộ qua dữ liệu, không qua chân reset
\end{itemize}

\paragraph{Hiện thực}\label{par:10-2}
\begin{itemize}
\item SystemVerilog tổng hợp được; FSM viết theo kiểu hai khối (\sig{always_ff} trạng thái, \sig{always_comb} trạng thái kế); không có latch. Cả bốn FSM mã hóa one-hot
\item \sig{job_t} khai báo trong \sig{gemm_pkg} là \sig{struct packed}; công cụ tổng hợp tự bỏ trường không dùng ở mỗi engine
\item Mọi so sánh idx dùng 32 bit không dấu
\item Không có bộ nhân trong FSM; đường tổ hợp dài nhất là \sig{want} → \sig{go} → trạng thái kế và \sig{blk_v}/\sig{blk_idx} (so sánh idx 32 bit), đủ cho chu kỳ 10 ns
\end{itemize}

\paragraph{Assertion bắt buộc (SVA)}\label{par:10-3}
\begin{spectable}{|L{0.1}|L{0.9}|}{Danh sách assertion của \sigc{gemm_main_fsm}}{tab:10-4}{\textbf{Mã} & \textbf{Khẳng định}}
AS-01 & \sig{o_rd_start} → kênh đọc rảnh; \sig{o_wr_start} → kênh ghi rảnh \\\hline
AS-02 & \sig{o_pe_start} → \sig{!i_pe_busy}; mỗi \sig{o_pe_start} có đúng một \sig{i_pe_done} \\\hline
AS-03 & Không bộ đếm \sig{rcnt} hay \sig{wcnt} nào tràn hay âm \\\hline
AS-04 & LOAD\_X bắt đầu vào bank b → không job nào ở slot, C\_RUN hay C\_HAND chưa tính xong có bank X = b (trừ job reuse; với X = BIAS chỉ xét job có \sig{post} = 1, vì job \sig{post} = 0 mang \sig{bbias} = 0 nhưng không đọc BIAS, D-35) \\\hline
AS-05 & COMPUTE vào bank C b bắt đầu → \sig{wcnt[C][b]} = 0 \\\hline
AS-06 & STORE\_C bắt đầu theo thứ tự idx tăng dần 1 \\\hline
AS-07 & \sig{o_done_pulse} ∧ \sig{o_err_pulse} không bao giờ cùng lúc; mỗi \sig{i_start_pulse} có tối đa một trong hai \\\hline
AS-08 & \sig{o_comp_*} ổn định trong khi C\_RUN; \sig{o_comp_dim[31:24]} = 0 \\\hline
AS-09 & Mỗi máy trạng thái (T\_*, R\_*, C\_*, S\_*) luôn là one-hot (\sig{onehot(state)}) \\\hline
AS-10 & \sig{o_err_pulse} → \sig{o_err_code} ∈ \{1, 2, 3, 4\} \\\hline
AS-11 & \sig{o_soft_clr} chỉ khi \sig{i_rd_busy} = 0 và \sig{i_wr_busy} = 0, và kéo dài đúng 1 chu kỳ \\\hline
AS-12 & T\_SOFT hoặc T\_CLR → không \sig{o_rd_start}, \sig{o_wr_start}, \sig{o_pe_start} \\\hline
AS-13 & \sig{blk_v} → không pha nào của job có idx ≥ \sig{blk_idx} bắt đầu \\\hline
AS-14 & \sig{o_done_pulse} → \sig{err_v} = 0 và job \sig{last} đã STORE\_C xong OKAY \\\hline
AS-15 & T\_IDLE → engine Đọc ở R\_IDLE ∧ engine Tính ở C\_IDLE ∧ engine Ghi ở S\_IDLE ∧ ¬\sig{slot_v} (D-34) \\\hline
\end{spectable}

\paragraph{Kiểm tra chức năng}\label{par:10-5}
\begin{itemize}
\item Chuỗi 1 … 7 descriptor ngẫu nhiên (stride khác 0, \sig{REUSE_B}, \sig{WAIT_PREV}, \sig{POST_EN}/\sig{RELU_EN}): DRAM sau DONE khớp thực thi tuần tự từng descriptor
\item Descriptor sai ở vị trí f: \sig{ERR_CODE} = 1, \sig{DESC_DONE_CNT} = f, \sig{ERR_DESC_ADDR} = địa chỉ descriptor f, không pha nào của job ≥ f bắt đầu sau lỗi, các job < f chạy xong
\item \sig{DESC_PTR} lệch căn 64 byte: \sig{ERR_CODE} = 1, không giao dịch AXI nào, \sig{o_err_pulse} ở t + 2
\item Lỗi AXI tiêm ở FETCH, LOAD\_*, STORE\_C cho mã 4, 2, 3 tương ứng
\item Soft reset ở T\_IDLE, ở T\_RUN khi kênh đọc bận, khi kênh ghi bận, khi COMPUTE đang chạy: sau \sig{o_soft_clr}, STATUS = 0, \sig{DESC_PTR} giữ nguyên, không DONE, không ERROR
\item Hai tác vụ liên tiếp: tác vụ 1 kết thúc ERROR khi slot hoặc C\_HAND còn giữ job idx ≥ f, tác vụ 2 chạy chuỗi hợp lệ: tác vụ 2 DONE, không có ghi nào vào \sig{C_ADDR} của tác vụ 1 sau lỗi; tương tự sau tác vụ DONE (D-34)
\item Mô hình COMPUTE: \sig{i_pe_done} cách \sig{o_pe_start} đúng COMPUTE\_CYCLES (534 chu kỳ với 32³)
\end{itemize}

\Needspace{10\baselineskip}
\begin{notebox}[Kiểm chứng bằng mô hình]
(Mô hình nằm ở \sig{gemm-accel/model/}, đã chạy lại với baseline ngày 12/11/2026; kết quả đạt.) \sig{gemm_sched_model.py} chạy 1 200 chuỗi ngẫu nhiên 1 – 7 descriptor (stride, \sig{REUSE_B}, \sig{WAIT_PREV} với descriptor đọc kết quả của descriptor trước, descriptor sai, lỗi AXI tiêm ngẫu nhiên ở FETCH/LOAD/STORE). Mọi bên đọc lấy dữ liệu ở đầu và cuối pha và phải thấy giống nhau, nên ghi đè sai bank bị phát hiện. Kết quả: DRAM khớp thực thi tuần tự ở mọi chuỗi DONE; với lỗi, \sig{DESC_DONE_CNT} = f, \sig{ERR_DESC_ADDR} đúng, không pha nào của job ≥ f bắt đầu sau lỗi. Mô hình đã bắt hai lỗi quy tắc trong quá trình thiết kế (đường \sig{REUSE_B} và lỗi cùng chu kỳ), dẫn tới bước R\_RB và tín hiệu \sig{blk_v}, \sig{blk_idx} tổ hợp.
\end{notebox}

\end{document}
